\section{The exact bridge between metric conditioning and observation balls}
\label{sec:modulus-bridge}
We now relate the three condition objects used in the paper. The
certificate $\eta_d$ is observable and controls all closed-model
competitors. The matrix $V_H$ is invariant under coefficient-chart
changes. The function $\omega_P$ is an exact target diameter. These
are three precise meanings, respectively, of representation-free,
coordinate-invariant, and decision-theoretic conditioning.

At a point of $\Theta_d^\circ$, partition the $kd$ locally labelled
roots into their equal-base-value blocks $B$. Let
$\Gamma=\prod_B\mathfrak S_B$, acting only within those blocks, and
write
\[
 d_\Gamma(x,y)=\min_{\pi\in\Gamma}\|x-\pi y\|_\infty,
 \qquad
 \mathcal E_V=\{V^{1/2}z:\|z\|_2\le1\},\qquad
 D_\Gamma(V)=\max_{x,y\in\mathcal E_V}d_\Gamma(x,y).
\]
The metric compares two perturbed targets, rather than one target
with the base point. This distinction is essential at collisions.

\begin{theorem}[Exact regular observation-ball modulus]
\label{thm:regular-modulus-bridge}
Let $P=F(\theta_0)$, $\theta_0\in\Theta_d^\circ$, and let
$H=H_{P,w}=\operatorname{diag}(w_j/P_{je})$ be the categorical Fisher
form. With the Hellinger normalization of
Section~\ref{sec:intrinsic-modulus},
\begin{equation}\label{eq:regular-modulus-constant}
 \omega_P(t)=2D_\Gamma(V_H(P))\,t+o(t).
\end{equation}
In particular, writing $\mathfrak c_H=\max_a\sqrt{(V_H)_{aa}}$,
\[
 2\mathfrak c_H\le
       \lim_{t\downarrow0}\frac{\omega_P(t)}t
       \le4\mathfrak c_H.
\]
If all aggregate base roots are distinct, the limit is exactly
$4\mathfrak c_H$. At a collision the constant is the stated finite
quotient functional, not an assumed multiple of a radial condition.
\end{theorem}
\begin{proof}
Theorem~\ref{thm:quotient-chart} accounts for the whole closed-model
observation neighbourhood. Write $I=J^{\mathsf T}HJ$. Uniformly in
a sufficiently small coefficient chart,
\[
 h_w(F(\theta_0+u),P)
       =\tfrac12\|Ju\|_H+O(\|u\|^2),\qquad
 g(\theta_0+u)=g(\theta_0)+Gu+O(\|u\|^2).
\]
The first expansion follows by Taylor expansion of every square root;
all base cells are positive. Injectivity of $J$ implies $\|u\|=O(t)$
for points in the Hellinger ball. Thus the rescaled coefficient balls
$t^{-1}(R(\mathcal B_P(t))-\theta_0)$ converge in Hausdorff distance
to $\{u:u^{\mathsf T}Iu\le4\}$. For the inner inclusion, first use
a strict ellipsoid and then let its radius increase to $2$; the outer
inclusion follows from the uniform expansion and the lower singular
value of $J$.

The image of $\{u:u^{\mathsf T}Iu\le1\}$ under $G$ is precisely
$\mathcal E_{V_H}$: the minimum of $u^{\mathsf T}Iu$ subject to
$Gu=x$ is $x^{\mathsf T}V_H^{-1}x$. Distinct base-value blocks have a
fixed positive gap. Hence for two sufficiently close representatives,
matching takes place within the blocks, and their target distance is
$d_\Gamma(Gu,Gv)+O(\|u\|^2+\|v\|^2)$. The quotient metric is
Lipschitz in both arguments. Hausdorff convergence and compactness
therefore give \eqref{eq:regular-modulus-constant}. Since $0$ belongs
to the centrally symmetric ellipsoid, its quotient diameter is at
least its sup-norm radius $\mathfrak c_H$, and the triangle inequality
gives at most twice that radius. Without collisions the pair $x,-x$
in a maximal-radius direction attains the latter bound.
\end{proof}

The constant in this theorem can be evaluated by finitely many convex
programs using only $V_H$. If $B$ has size $b$, let $x_{B,(i)}$ denote
its $i$th increasing coordinate and put
\[
 M_{B,i}=\max_{x\in\mathcal E_{V_H}}x_{B,(i)}.
\]
Then
\begin{equation}\label{eq:finite-quotient-constant}
 D_\Gamma(V_H)=\max_{B,\,1\le i\le b}
                     \{M_{B,i}+M_{B,b+1-i}\}.
\end{equation}
Indeed sorting is an isometry for the real bottleneck metric, and
central symmetry changes the minimum of the $i$th order statistic
into the negative maximum of the $(b+1-i)$th. More explicitly,
$M_{B,i}$ is the maximum, over subsets $S\subset B$ of cardinality
$b-i+1$, of
\[
 \max\{a:\|z\|_2\le1,\ (V_H^{1/2}z)_r\ge a\ (r\in S)\}.
\]
Each inner problem is a second-order cone problem. This is a finite
observable formula, not quantifier elimination over latent factors.
For example, for one two-coordinate collision block and
$V=\operatorname{diag}(1,\epsilon^2)$, $0<\epsilon\le1$,
\[
 D_\Gamma(V)=1+\frac{\epsilon}{\sqrt{1+\epsilon^2}}<2.
\]
This example concerns the covariance functional itself and explains
why a universal collision-blind factor $4\mathfrak c_H$ would be
incorrect.

\begin{proposition}[Global certificate comparison and its limitation]
\label{prop:certificate-modulus}
At a full-rank true datum with $\tau(P)>0$, the global observable
theorem implies, for every $t\ge0$,
\begin{equation}\label{eq:eta-to-modulus}
 \omega_P(t)\le C\min\{1,(t/\eta_d(P))^{1/d}\}.
\end{equation}
On the within-component $\nu$-separated class the exponent is $1$;
the cluster versions have exponent $1/m$ under their stated margins.
Consequently $2D_\Gamma(V_{H_{P,w}})\le C_\nu/\eta_d(P)$ on the
regular stratum. There is no uniform reverse comparison of
$\eta_d(P)^{-1}$ with this Fisher-metric constant over all strictly
positive regular data, even when the component roots are fixed.
\end{proposition}
\begin{proof}
For probability vectors $p,q$, factoring $p_e-q_e$ into the difference
and sum of their square roots gives $\|p-q\|_2\le2h(p,q)$. Hence a
member of $\mathcal B_P(t)$ has observation error at most
$2t/\sqrt{\min_jw_j}$ from $P$. Apply the global observable estimate
from $P$ to each member of a pair and use the triangle inequality.
The separated and clustered versions follow from the same argument.
The regular limit is Theorem~\ref{thm:regular-modulus-bridge}.

For the last assertion fix an interior regular binary, two-component,
degree-one datum $\overline P$ with distinct roots, and replace its
first channel by $D_\epsilon\overline U$, where
\[
 D_\epsilon=\begin{pmatrix}1&1-\epsilon\\0&\epsilon\end{pmatrix},
 \qquad P_\epsilon=D_\epsilon\overline P,\quad 0<\epsilon<1.
\]
Both channels remain strictly positive and full rank; the roots and
weights are unchanged. The normalization residual has the same
kernel because $D_\epsilon$ is invertible. In square coordinates the
leading matrix is $A_\epsilon=D_\epsilon\overline A$, so
$B_1(P_\epsilon)\ge\|A_\epsilon^{-1}\|_\op\ge c/\epsilon$.
Thus $\eta_1(P_\epsilon)^{-1}\ge c/\epsilon$.

The row transformation is a local diffeomorphism between the two
parameter charts. For an observation tangent $Z_\epsilon$,
$\overline Z=D_\epsilon^{-1}Z_\epsilon$ has second row
$Z_{\epsilon,2}/\epsilon$ and first row
$Z_{\epsilon,1}-(1-\epsilon)Z_{\epsilon,2}/\epsilon$.
The second row probabilities of $P_\epsilon$ are
$\epsilon\overline P_2$, whereas its first row probabilities have a
fixed positive lower bound. It follows, clock by clock, that
\[
 \|\overline Z\|_{H_{\overline P,w}}
        \le C\epsilon^{-1/2}\|Z_\epsilon\|_{H_{P_\epsilon,w}}.
\]
The target is unchanged by the row transformation, so
$\mathfrak c_{H_{P_\epsilon,w}}(P_\epsilon)
\le C\epsilon^{-1/2}$. The regular modulus constant is at most four
times this quantity, while $\eta_1^{-1}\ge c/\epsilon$.
Their ratio tends to zero. The failure occurs as a cell floor vanishes;
no reverse comparison on a fixed positive-cell compact class is
being disproved by this example.
\end{proof}
